\appendix
	
	\section{Firedrake Implementation}
	
The different formulations introduced in the previous sections have been implemented using the finite element platform Firedrake \cite{firedrake}. This implementation is significantly simplified by the availability of automatic differentiation tools. The approach involves defining the discrete Lagrangian in equation (\ref{eq:discLag}) and computing its variations using a \texttt{derivative} function.

For clarity, the implementation is organized into several sections. Some details are omitted for simplicity, but the complete code can be provided upon request.

	In the first place all the physical and numerical parameters must be properly
	initialize:
	\lstinputlisting[language=Python]{Text/probdef.py}
	The keyword \texttt{formulation} can also be set to \texttt{\textquotesingle Natural\textquotesingle} and the different parameters set to their desired 
	value, depending on the formulation adopted.
	The user must also specify two functions defining the fields
	$\etP$ and $\stP$ and the source term $q$ in the balance equation, 
	i.e.,
	\lstinputlisting[language=Python]{Text/eTsT.py}	
	The \texttt{ufl} language \cite{ufl_fenics} allows the
	expressive specification of the operators
	involved in the variational formulation.
	The first thing is to specify the finite element spaces,
	the symbols for the different variables $\stP, \uP, \eP, \lP$ and $\mP$ 
	and the expression of the Lagrangian function: 
	\lstinputlisting[language=Python]{Text/lagrangian.py}
	The boundary conditions must also be specified. In this simple
	example homogeneous Dirichlet boundary conditions are imposed on 
	the whole boundary, however more sophisticated 
	cases, like mixed Dirichlet-Neumann conditions 
	can easily be incorporated. 
	The problem can be finally solved by invoking
	a \texttt{PETSc} solver available in the \texttt{Firedrake}
	platform:
	\lstinputlisting[language=Python]{Text/solve.py}
	This formulation can be change to account for non-linear effects.
	This would be the case for instance if the reaction term above depends 
	in a non-linear manner of $u$.
	
	
